\chapter{Dodatne podrobnosti grafovske konstrukcije}
\label{app:grafi}

Ta dodatek dopolnjuje opis grafovske konstrukcije iz poglavja~\ref{sec:grafovska-predstavitev}.

\section{Razširjene povezave znotraj fiksacije}
\label{app:razsirjene-fiksacijske-povezave}

Fiksacijo obravnavamo kot neprekinjeno zaporedje vozlišč z istim veljavnim identifikatorjem fiksacije.
Če ima taka fiksacija $F$ vozlišč, jim priredimo lokalne indekse $0,\ldots,F-1$.
Parameter $k_f$ določa število fiksacijskih povezav, ki jih generira vsako vozlišče.
Za posamezno fiksacijo najprej izračunamo korak razširitve
\begin{equation}
    s = \max(1, \lfloor F/k_f + 0{,}5 \rfloor).
\end{equation}
Množico kandidatnih odmikov določimo kot
\begin{equation}
    D = \{(1 + q s) \bmod F \mid q = 0,\ldots,k_f-1\}.
\end{equation}
Morebitni odmik $0$ in podvojene odmike odstranimo.
Za vsak odmik $d \in D$ dodamo neusmerjene povezave
\begin{equation}
    \{v_i, v_{(i+d)\bmod F}\},
    \qquad i = 0,\ldots,F-1.
\end{equation}
Pri zapisu grafa za posredovanje sporočil vsako neusmerjeno povezavo predstavimo z dvema usmerjenima povezavama in odstranimo podvojene povezave znotraj relacije $E_f$.

Pri fiksaciji dolžine $F=21$ in parametru $k_f=3$ dobimo $s=7$.
Vozlišče z lokalnim indeksom $0$ se zato poveže z vozlišči $1$, $8$ in $15$.
Odmik $1$ ohrani časovno lokalno komunikacijo, večji odmiki pa omogočijo širjenje informacij tudi med časovno bolj oddaljenimi meritvami iste fiksacije.
Ker vsaka plast GNN omogoča izmenjavo informacij med neposredno povezanimi vozlišči, se doseg znotraj fiksacije dodatno povečuje s številom plasti.

\section{Ločeni prikazi tipov povezav}

Slike prikazujejo isti izsek časovnega okna, vendar ločeno za posamezne tipe povezav.
Namenjene so predvsem lažjemu branju slike~\ref{fig:konstrukcija-grafa-primer}: časovne povezave pokažejo lokalni red meritev, prostorske povezave sosednost položajev pogleda na zaslonu, fiksacijske povezave pa povezovanje meritev znotraj iste zaznane fiksacije.

% Vir slik: GFM-for-eyetracker/results/quick_v1_v2_comparison/RETAIN_2026-06-12_16-29-08/model_runs/all_signals/2026-06-12_18-09-29/experiments/multiclass_table6_valence_3class_emotion-elicitation/snapshot.csv.
% Slike so generirane s skriptom GFM-for-eyetracker/scripts/create_thesis_gnn_graph_examples.py.
\begin{figure}
    \centering
    \includegraphics[width=\textwidth]{gnn_graph_far_02_P10_90_avi_t106_86-107_95_kt1_ks1_kf3_f3_n9/casovne.png}
    \caption{Časovne povezave v izseku grafa. Oranžne povezave vodijo naprej po zaporedju meritev, rožnate pa nazaj, zato prikaz poudari lokalno časovno okolico vsakega vozlišča.}
    \label{fig:casovne-povezave-primer}
\end{figure}

\begin{figure}
    \centering
    \includegraphics[width=\textwidth]{gnn_graph_far_02_P10_90_avi_t106_86-107_95_kt1_ks1_kf3_f3_n9/prostorske.png}
    \caption{Prostorske povezave v izseku grafa. Povezujejo meritve s podobnim položajem pogleda na zaslonu, zato lahko povežejo tudi časovno bolj oddaljena vozlišča.}
    \label{fig:prostorske-povezave-primer}
\end{figure}

\begin{figure}
    \centering
    \includegraphics[width=\textwidth]{gnn_graph_far_02_P10_90_avi_t106_86-107_95_kt1_ks1_kf3_f3_n9/fiksacijske.png}
    \caption{Fiksacijske povezave v izseku grafa. Povezujejo meritve iz iste zaznane fiksacije, vendar z omejenim številom odmikov, da graf ne postane preveč gost.}
    \label{fig:fiksacijske-povezave-primer}
\end{figure}
